\documentclass[11pt]{article}
\usepackage[margin=1in]{geometry}
\usepackage{amsmath,amssymb,amsthm}
\usepackage[colorlinks=true,linkcolor=blue,urlcolor=blue,citecolor=blue]{hyperref}
\newtheorem{theorem}{Theorem}[section]
\newtheorem{lemma}[theorem]{Lemma}
\newtheorem{proposition}[theorem]{Proposition}
\newtheorem{corollary}[theorem]{Corollary}
\newtheorem{remark}[theorem]{Remark}
\numberwithin{equation}{section}
\newcommand{\tr}{\operatorname{tr}}
\newcommand{\rank}{\operatorname{rank}}
\newcommand{\sgn}{\operatorname{sgn}}
\newcommand{\sech}{\operatorname{sech}}
\newcommand{\op}{\mathrm{op}}
\newcommand{\supp}{\operatorname{supp}}
\newcommand{\C}{\mathbb C}
\newcommand{\R}{\mathbb R}
\newcommand{\cH}{\mathcal H}
\newcommand{\cN}{\mathcal N}
\newcommand{\cW}{\mathcal W}
\newcommand{\cA}{\mathcal A}
\newcommand{\diag}{\operatorname{diag}}
\newcommand{\Rea}{\operatorname{Re}}
\title{Analytic channels and exact constants\protect\\
for homogeneous fractional matrix maps}
\author{Working research manuscript}
\date{November 2026}
\begin{document}
\maketitle

\begin{abstract}
Fix $0<\alpha<1$ and a complex homogeneous function
$f(x)=a x_+^\alpha+b(-x)_+^\alpha$ on the real line, written as
$f=\kappa(-ix)^\alpha+\lambda(ix)^\alpha$. Uniformly over all nonzero
coefficient pairs, its optimal dimension-free Schatten power constant
satisfies
$s^2\cH_s(f)/(|\kappa|^s+|\lambda|^s)
\to1/[\pi^2\alpha(1-\alpha)]$ as $s\downarrow0$.
For fixed functions this gives exactly two exceptional complex lines
with half the generic coefficient. We resolve the transition as a
branch vanishes, determine sharp approximation ranks, and prove the
joint law when the matrix dimension is comparable to $s^{-2}$.
The lower proof retains the nonlinear rank-one resolvent denominator
and uses classical Cauchy-kernel localization to separate its periodic
singular channels. A positive rational Gaussian rule realizes any
$q$-pole rank-one update quadrature's singular spectrum in at most $q$ states,
including conjugate poles, and proves the dimension transition.
A bounded-input semifinite extension identifies the role of matrix
degree. The Gaussian localization input and PSD comparison theorem
are stated precisely from a companion manuscript. Classical scalar
approximation and rational quadrature methods are antecedents.
\end{abstract}

\section{Statements and context}

Throughout, $\alpha\in(0,1)$ and $\beta=1-\alpha$ are fixed. For a
matrix $T$, write $|T|=(T^*T)^{1/2}$ and
$\|T\|_s=(\tr|T|^s)^{1/s}$ for $s>0$. Let
\begin{equation}\label{eq:f}
 f(x)=a x_+^\alpha+b(-x)_+^\alpha,\qquad a,b\in\C.
\end{equation}
We distinguish fixed-function limits from limits allowing coefficients
to vary. Functional calculus is
applied to Hermitian inputs; the outputs need not be Hermitian. Define
\begin{equation}\label{eq:H}
 \cH_s(f)=\sup_{n\ge1}\sup_{\substack{A=A^*,\ B=B^*\\A\ne B}}
 \frac{\tr|f(B)-f(A)|^s}{\tr|B-A|^{\alpha s}},
\end{equation}
and let $\cH_{s,n}(f)$ restrict the supremum to dimension at most $n$.
Zero padding makes this equivalent to dimension exactly $n$.
The usual optimal constant in
$\|f(B)-f(A)\|_s\le K_s(f)\|B-A\|_{\alpha s}^\alpha$ is
$K_s(f)=\cH_s(f)^{1/s}$; it is not the power constant in \eqref{eq:H}.

Put $\phi=\pi\alpha/2$ and define the two principal branches
\[
 g_+(x)=(-ix)^\alpha
 =\cos\phi\,|x|^\alpha-i\sin\phi\,\sgn(x)|x|^\alpha,
 \qquad g_-(x)=\overline{g_+(x)}.
\]
They vanish at zero. With $e=(a+b)/2$ and $o=(a-b)/2$, set
\begin{equation}\label{eq:channels}
 \kappa=\frac12\left(\frac e{\cos\phi}
                         +\frac{io}{\sin\phi}\right),\qquad
 \lambda=\frac12\left(\frac e{\cos\phi}
                         -\frac{io}{\sin\phi}\right),\qquad
 m=1_{\kappa\ne0}+1_{\lambda\ne0}.
\end{equation}
Then $f=\kappa g_++\lambda g_-$. The case $m=0$ is exactly $f=0$.
For nonzero $f$, $m=1$ precisely on the two lines
$b=a e^{i\pi\alpha}$ and $b=a e^{-i\pi\alpha}$; otherwise $m=2$.
Write
\begin{equation}\label{eq:cf}
 b_\alpha=\frac1{2\pi^2\alpha\beta},\qquad c_f=m b_\alpha.
\end{equation}

For a positive integer $r$, let $\cN_{f,r,n}(L)$ be the supremum of
$\#\{j:\sigma_j(f(B)-f(A))>e^{-L}\}$ over Hermitian pairs of size at
most $n$ satisfying $\rank(B-A)\le r$ and $\|B-A\|_\op=1$.
Omit $n$ for unrestricted finite dimension.

\begin{theorem}[Uniform coefficient law]\label{thm:uniform}
For fixed $\alpha\in(0,1)$,
\begin{equation}\label{eq:uniform}
 \lim_{s\downarrow0}\ \sup_{(\kappa,\lambda)\ne(0,0)}
 \left|\frac{s^2\cH_s(\kappa g_++\lambda g_-)}
 {|\kappa|^s+|\lambda|^s}-2b_\alpha\right|=0.
\end{equation}
The coefficients may depend arbitrarily on $s$; no uniformity in
$\alpha$ is asserted.
\end{theorem}

For $|\theta|\le1$, put $f_\theta=g_++\theta g_-$, with
$\log0=-\infty$. Every nonzero function reduces to this form by scalar
homogeneity and exchanging $g_+,g_-$ using $x\mapsto-x$.
The uniform accuracy law retains the smaller amplitude explicitly.

\begin{theorem}[Accuracy transition]\label{thm:accuracy}
For every fixed positive integer $r$,
\begin{equation}\label{eq:accuracy}
 \lim_{L\to\infty}\ \sup_{|\theta|\le1}
 \left|\frac{\cN_{f_\theta,r}(L)}{L^2}
 -r b_\alpha\left[1+\left(1+\frac{\log|\theta|}{L}\right)_+^2\right]
 \right|=0.
\end{equation}
In particular, if $-L^{-1}\log|\theta_L|\to\tau\in[0,\infty]$,
the count coefficient is $r b_\alpha[1+(1-\tau)_+^2]$.
At $\theta=0$ or $\tau=\infty$ the second term is zero.
All coefficient phases are allowed.
\end{theorem}

\begin{corollary}[Fixed functions]\label{thm:main}
For \eqref{eq:f} with fixed $\alpha,a,b$,
\begin{equation}\label{eq:main}
 \lim_{s\downarrow0}s^2\cH_s(f)=2c_f,\qquad
 \lim_{L\to\infty}\frac{\cN_{f,r}(L)}{L^2}=r c_f
 \quad\text{for each fixed positive integer }r.
\end{equation}
\end{corollary}

Every nonzero \emph{real}-valued function \eqref{eq:f}, including
$(x_+)^\alpha$, has $m=2$. For comparison, the
positive-input power theorem in the companion paper \cite{Companion}
has moment coefficient $1/(2\pi^2\alpha\beta)$ and count coefficient
$1/(4\pi^2\alpha\beta)$. Thus the moment coefficients for PSD powers,
one analytic branch, and all nonzero real homogeneous maps have ratio
$1:2:4$. This is a comparison of power constants. For nonzero $f$,
\begin{equation}\label{eq:K}
 \log K_s(f)=\frac2s\log\frac1s+\frac{\log(2c_f)}s+o(s^{-1}).
\end{equation}

For $u\in[0,\infty]$ and $x\ge0$, define
\begin{equation}\label{eq:pu}
 p_u(x)=b_\alpha[x^2+(x-u)_+^2],\qquad
 p_\infty(x)=b_\alpha x^2,\qquad
 F(t,u)=\int_0^\infty e^{-x}\min\{t,p_u(x)\}\,dx.
\end{equation}

\begin{theorem}[Joint dimension and coefficient transition]\label{thm:joint}
Suppose $s\downarrow0$, $n_s s^2\to t\in(0,\infty)$,
$|\theta_s|\le1$, and $u_s=-s\log|\theta_s|\to u\in[0,\infty]$.
Then
\begin{equation}\label{eq:joint}
 s^2\cH_{s,n_s}(f_{\theta_s})\longrightarrow F(t,u).
\end{equation}
If $a_0>0$ is the solution of $p_u(a_0)=t$, the formula is
\begin{align}
 F(t,u)&=2b_\alpha[1-(1+a_0)e^{-a_0}],
 &a_0&=\sqrt{t/b_\alpha}, && t\le b_\alpha u^2,\label{eq:jointone}\\
 F(t,u)&=2b_\alpha[1+e^{-u}-(2+2a_0-u)e^{-a_0}],
 &a_0&=\frac{u+\sqrt{2t/b_\alpha-u^2}}2,
 &&t>b_\alpha u^2.\label{eq:jointtwo}
\end{align}
For $u=\infty$ use \eqref{eq:jointone} for all finite $t$.
For every $0<t_0<t_1<\infty$,
\begin{equation}\label{eq:jointuniform}
 \sup_{\substack{|\theta|\le1,\ n\in\mathbb N\\t_0\le n s^2\le t_1}}
 \left|s^2\cH_{s,n}(f_\theta)
          -F(n s^2,-s\log|\theta|)\right|\longrightarrow0.
\end{equation}
If instead $n_s s^2\to\infty$, the limit is
$2b_\alpha(1+e^{-u})$.
\end{theorem}

For arbitrary coefficient sequences, let
\[
 M_s=\max(|\kappa_s|,|\lambda_s|)>0,\qquad
 u_s=-s\log\frac{\min(|\kappa_s|,|\lambda_s|)}{M_s}.
\]
After exchanging branches and multiplying by a scalar phase,
\eqref{eq:joint} means
\begin{equation}\label{eq:jointscale}
 s^2\cH_{s,n_s}(\kappa_s g_++\lambda_s g_-)
       =M_s^s\,[F(t,u)+o(1)].
\end{equation}
The error is normalized by $M_s^s$; it is not an absolute $o(1)$
when that factor is unbounded. The case of a zero smaller coefficient
means $u_s=\infty$.

\begin{corollary}[Fixed-function dimension limits]\label{thm:dimension}
Suppose $f\ne0$ is fixed, $s\downarrow0$, and
$n_s s^2\to t\in(0,\infty)$.
Then
\begin{equation}\label{eq:transition}
 s^2\cH_{s,n_s}(f)\longrightarrow
 \Phi_{c_f}(t),\qquad
 \Phi_c(t)=2c\left[1-\left(1+\sqrt{t/c}\right)e^{-\sqrt{t/c}}\right].
\end{equation}
If $n_s s^2\to0$, then $\cH_{s,n_s}(f)/n_s\to1$; if
$n_s s^2\to\infty$, then $s^2\cH_{s,n_s}(f)\to2c_f$.
For fixed $r$ and $n_L/L^2\to t\in(0,\infty)$,
\begin{equation}\label{eq:cappedcount}
 \frac{\cN_{f,r,n_L}(L)}{L^2}\longrightarrow\min(t,r c_f).
\end{equation}
\end{corollary}

The counts are the optimal universal ranks for operator error
$\varepsilon\|B-A\|_\op^\alpha$, with $\varepsilon=e^{-L}$, by the
best-rank approximation theorem. They do not measure relative error
against the possibly much smaller correction norm.

\paragraph{Antecedents and scope.}
Dimension-free H\"older estimates in quasi-Banach Schatten classes
are established work \cite{Ricard,HSZ}. Root-exponential fractional
correction estimates already occur in Birman--Koplienko--Solomyak
\cite{BKS}. For the kernel $x^p y^q/(x+y)^\nu$ on $(0,1)^2$,
Laptev \cite{Laptev} proves
$N(\varepsilon)\sim[2\pi^2\varrho]^{-1}\log^2(1/\varepsilon)$
when $\nu>0$, $p,q>-1/2$, and $\varrho=p+q-\nu+1>0$.
His proof uses logarithmic hyperbolic-cosine reduction,
diagonal-cell localization, and lower-order off-diagonal terms.
The localization mechanism used below belongs to
that classical spectral-asymptotic framework; we do not claim a new
Weyl or localization method. Sharp scalar rational approximation is
also classical:
Pekarskii \cite{Pekarskii} treats endpoint-weighted Markov functions,
Stahl \cite{Stahl} treats powers, and Mardvilko \cite{Mardvilko}
proves matching weak rates for odd and even power extensions through
Cauchy transforms. These scalar rates are not asserted as new here.
Rank-one functional-calculus differences and two Hankel/Riesz
channels also have established connections; see Arazy--Zelenko
\cite[Introduction and Section 4]{AZ}. Rational matrix-update exactness
\cite{BCKS} and positive rational
Gaussian quadrature \cite{Gautschi} are further ingredients with prior
literature. The results here concern exact universal perturbation
constants, their analytic exceptional lines, and joint dimension limits.
The uniform formulas resolve approach to an exceptional line on the
scales $-s\log|\theta|$ and $-L^{-1}\log|\theta|$.
The proofs have received internal mathematical reviews; no priority
certificate or external peer-review claim is made. Shifted solves must
be paid for separately, and no finite-word stability or general quadratic
matrix algorithm is asserted.

\section{Uniform quadrature and the upper bounds}

We use the following precise localization statement from
\cite[Lemma 3.1]{Companion}; only this quadrature theorem is imported.
\begin{lemma}[Gaussian localization]\label{lem:gaussian}
Let a Banach-valued $F$ be analytic on $|\Im z|<d_0$. Suppose that for
every fixed $0<d<d_0$,
\[
 \|F(x+i\eta)\|\le M_d e^{-u(-x)_+-v x_+},
 \qquad |\eta|\le d,\quad u,v>0.
\]
For all sufficiently large $L$ there is a quadrature with positive
weights and real nodes, of error at most $e^{-L}$, using
\begin{equation}\label{eq:qnodes}
 q_L\le\frac{L^2}{4\pi d}\left(\frac1u+\frac1v\right)
                    +O_{u,v,d}(L^{7/4})
\end{equation}
nodes. The threshold may depend on $M_d$, uniformly for a common
envelope. Nodes and weights can be chosen using only the envelope
parameters and $L$, as in the cited construction; this same quadrature
then works for every integrand with that envelope.
Each node of rank at most $k$ gives a quadrature of rank
at most $kq_L$.
\end{lemma}

Here is its application, including the needed envelope. With
$R_A(z)=(A-zI)^{-1}$ and $c_\alpha=\sin(\pi\alpha)/\pi$,
\begin{equation}\label{eq:onepole}
 g_+(B)-g_+(A)=-i c_\alpha\int_0^\infty t^\alpha
                  [R_B(-it)-R_A(-it)]\,dt.
\end{equation}
For a scalar $x$ this follows by applying the fractional-power integral
to $-ix$; equivalently the even and odd beta integrals combine and
cancel the $+it$ pole. Functional calculus gives the matrix identity.
The $g_-$ formula is its adjoint.

Normalize $\|B-A\|_\op\le1$. If $t=\rho e^{i\eta}$ and
$|\eta|\le d<\pi/2$, then
$\|R_A(\pm it)\|\le(\rho\cos d)^{-1}$. Bounding two resolvents
separately, and then using their difference identity, gives
\[
 \|R_B(\pm it)-R_A(\pm it)\|
 \le C_d\min(\rho^{-1},\rho^{-2}).
\]
After $t=e^z$, the combined integrand for $f$ is analytic on
$|\Im z|<\pi/2$ and has bound
\[
 C_{f,d}\min\{e^{\alpha x},e^{-\beta x}\}.
\]
This also supplies a fixed output norm bound $M_f\ge1$.
Each active pole difference has rank at most $\rank(B-A)$.
Lemma~\ref{lem:gaussian} therefore yields, for fixed $d<\pi/2$,
\begin{equation}\label{eq:uppercount}
 N_\Delta(L)\le r c_d L^2+O_{f,d,r}(L^{7/4})+O_{f,d,r}(1),
 \qquad c_d=\frac{m}{4\pi d\alpha\beta},
\end{equation}
uniformly in all normalized rank-$r$ pairs. A fixed constant enlargement
extends it to all $L\ge0$. Notice $c_d\downarrow c_f$ at the strip
boundary; that boundary is approached only after the large-$L$ limit.

We also need an upper uniform in a vanishing branch. Put
\begin{equation}\label{eq:bd}
 b_d=\frac1{4\pi d\alpha\beta},\qquad 0<d<\pi/2.
\end{equation}
For $f_\theta=g_++\theta g_-$, $|\theta|\le1$, use separate
quadratures with depths
\[
 h_+=L+C_d,\qquad h_-=[L+\log|\theta|]_++C_d.
\]
Choose the fixed $C_d$ to pass the Gaussian threshold and make the
sum of the weighted errors at most $e^{-L}$. Indeed
$|\theta|e^{-h_-}\le e^{-L-C_d}$ even when
$L+\log|\theta|<0$. For $\theta=0$ omit the second rule.
The node count and the correction rank therefore give
\begin{equation}\label{eq:uniformupper}
 \cN_{f_\theta,r}(L)\le
 r b_d\{L^2+[L+\log|\theta|]_+^2\}
       +O_{\alpha,d,r}(L^{7/4}+1),\qquad L\ge0,
\end{equation}
uniformly in $\theta$ and the background. The two branch envelopes
also give a common output norm bound $M_\alpha\ge1$ for all these
unit-input corrections.

For $0<s\le1$, the power triangle
$\tr|X+Y|^s\le\tr|X|^s+\tr|Y|^s$ and homogeneity give the exact
rank-one reduction
\begin{equation}\label{eq:rankone}
 \cH_s(f)=\sup_{A=A^*,\ \|v\|=1}
       \tr|f(A+vv^*)-f(A)|^s.
\end{equation}
Indeed telescope the signed spectral decomposition of $B-A$; a negative
increment is a reversed positive one, and its contribution scales by
the $\alpha s$ power of its size. The sum of those powers is the
denominator in \eqref{eq:H}. This works within each fixed dimension.

For a correction of norm at most one, layer cake reads
\begin{equation}\label{eq:layer}
 \tr|\Delta|^s=s\int_0^\infty e^{-sL}N_\Delta(L)\,dL.
\end{equation}
For norm bound $M_f>1$ divide by $M_f$, which shifts logarithmic
thresholds by a constant and multiplies the moment by $M_f^s\to1$.
Integrating \eqref{eq:uppercount} for rank one gives
\[
 \cH_s(f)\le\frac{2c_d}{s^2}+O_{f,d}(s^{-7/4})+o(s^{-2}).
\]
Take $s\downarrow0$ before $d\uparrow\pi/2$. This proves both upper
bounds in Corollary~\ref{thm:main}; $f=0$ is immediate.
Finally, the power triangle and the adjoint relation of the branches
give
\begin{equation}\label{eq:coefficientupper}
 \cH_s(\kappa g_++\lambda g_-)
 \le (|\kappa|^s+|\lambda|^s)\cH_s(g_+).
\end{equation}
The one-branch upper just proved has
$\limsup s^2\cH_s(g_+)\le2b_\alpha$. This proves the uniform upper
in Theorem~\ref{thm:uniform}. Equation~\eqref{eq:uniformupper}, followed
by $d\uparrow\pi/2$, proves the uniform upper in
Theorem~\ref{thm:accuracy}.

\section{An exact nonlinear centered factorization}

Fix $\delta>0$ and the probability density
\begin{equation}\label{eq:rho}
 \rho(x)=\begin{cases}c_0e^x,&x\le0,\\c_0e^{-\delta x},&x\ge0,
 \end{cases}\qquad c_0=(1+\delta^{-1})^{-1}.
\end{equation}
On a bounded interval $I$, let $D$ multiply by $e^x$ on $L_2(I)$,
$w=\sqrt\rho|_I$, and set
\begin{equation}\label{eq:pair}
 H=\begin{pmatrix}0&D\\D&0\end{pmatrix},\quad
 v=(w,0),\quad E=vv^*,\quad A=H-E/2,\quad B=H+E/2.
\end{equation}
The rank-one increment has norm $\int_I\rho\le1$.
Define the \emph{actual} nonlinear weights
\begin{equation}\label{eq:omega}
 M_I(t)=t\int_I\frac{\rho(x)}{e^{2x}+t^2}\,dx,
 \qquad \omega_I(t)=(1+M_I(t)^2/4)^{-1}.
\end{equation}
They obey $\omega_0\le\omega_I\le1$ with a positive constant independent
of $I,t$: increasing $I$ and bounding $e^{-\delta x}\le1$ on the
right gives
\[
 M_I(t)\le c_0\left[\frac\pi2+
                   \frac{\log(1+t^2)}{2t}\right]\le\frac\pi2+1.
\]

Write $R(t)=(D^2+t^2)^{-1}$ and $J=\diag(-I,I)$. Centered Woodbury
subtraction gives
\[
 R_B(z)-R_A(z)=-\frac{R_H(z)ER_H(z)}{1-(v^*R_H(z)v)^2/4}.
\]
At $z=-it$, $v^*R_H(-it)v=-iM_I(t)$ and
$R_H(-it)v=(-itR(t)w,DR(t)w)$. Inserting this in
\eqref{eq:onepole} proves
\begin{equation}\label{eq:RJ}
 \Delta g_+=i\mathcal R J,\qquad \mathcal R\succeq0,
 \qquad \Delta g_-=-iJ\mathcal R.
\end{equation}
Thus $\Delta f$ has the singular values of
$\kappa\mathcal R-\lambda J\mathcal R J$.

For the explicit factors put
\begin{align}
 u(x)&=\sqrt{c_\alpha\rho(x)}e^{-x/2},&
 v_I(y)&=e^{\alpha y/2}\sqrt{\omega_I(e^y)},\label{eq:weights}\\
 k_+(z)&=\frac{e^{-z/2}}{2\cosh z},&
 k_-(z)&=\frac{e^{z/2}}{2\cosh z}.\nonumber
\end{align}
Define $T_\pm:L_2(\R)\to L_2(I)$ by their kernels
$u(x)k_\pm(x-y)v_I(y)$. The formal notation $M_u C_\pm M_{v_I}$
always means these bounded kernel extensions; $M_{v_I}$ alone is
unbounded on the whole line. They are Hilbert--Schmidt uniformly in
$I$. Indeed
\begin{equation}\label{eq:HS}
 \|T_\pm\|_{\mathrm{HS}}^2\le
 \left(\int_\R u(x)^2e^{\alpha x}\,dx\right)
 \left(\int_\R e^{-\alpha z}|k_\pm(z)|^2\,dz\right)<\infty.
\end{equation}
The first integral has slopes $\alpha$ and $\beta+\delta$; the
second is finite because the slowest one-sided kernel decay is
$e^{-|z|/2}$ and $\alpha<1$.

Substitution $t=e^y$ in \eqref{eq:RJ} gives
\begin{align}\label{eq:exactfactor}
 \mathcal R=VWW^*V^*,\quad W=\begin{bmatrix}T_+\\T_-\end{bmatrix},
 \quad V=\diag(I,iI),\nonumber\\
 \cW=[W,JW]=\diag(T_+,T_-)H_0,
 \quad H_0=\begin{pmatrix}1&-1\\1&1\end{pmatrix}.
\end{align}
For example, the two real kernels before the factor $V$ have numerators
$\sqrt{c_\alpha\rho(x)}e^{(\alpha+3)y/2}$ and
$\sqrt{c_\alpha\rho(x)}e^{x+(\alpha+1)y/2}$, divided by
$e^{2x}+e^{2y}$ and multiplied by $\sqrt{\omega_I(e^y)}$.
Since $V$ commutes with $J$, the singular values of the correction are
exactly those of
\begin{equation}\label{eq:pencil}
 \cW K\cW^*,\qquad K=\diag(\kappa,-\lambda).
\end{equation}
All these identities retain \eqref{eq:omega}; no linearization or
deletion of a weight inside an indefinite expression has occurred.

\section{Classical Cauchy localization with explicit rank bounds}

We give the needed localization estimates explicitly, within the
classical framework of \cite{Laptev}, to track the actual nonlinear
weights and the matrix channels. On $z\ge0$,
\begin{equation}\label{eq:expkernel}
 k_+(z)=\frac{e^{-3z/2}}{1+e^{-2z}},\qquad
 k_-(z)=\frac{e^{-z/2}}{1+e^{-2z}};
\end{equation}
on $z\le0$ use reflection and interchange the signs. The polynomial
\begin{equation}\label{eq:Pn}
 P_N(t)=\frac23\sum_{j=0}^N\left(-\frac{2t-1}{3}\right)^j
 \quad\text{satisfies}\quad
 \sup_{0\le t\le1}\left|\frac1{1+t}-P_N(t)\right|\le3^{-N-1}.
\end{equation}
Consequently each fixed-sign kernel is approximated by $N+1$
separated exponentials, with error at most $3^{-N-1}e^{-\gamma|z|}$,
where $\gamma$ is $1/2$ or $3/2$.

\begin{lemma}\label{lem:local}
Let $I$ be a union of $J_0$ consecutive cells of common length
$\ell\ge1$. On each cell periodize $k_\pm$ with period $\ell$,
and let $A_\pm$ be the resulting weighted block diagonal operators,
extended by zero on inputs outside $I$. Then
\begin{equation}\label{eq:local}
 T_\pm=A_\pm+F_\pm+E_\pm,\qquad
 \rank F_\pm\le4J_0(N+1),\qquad
 \|E_\pm\|\le C\sqrt\ell\,e^{\alpha\ell/2}3^{-N},
\end{equation}
where $C$ is independent of $I,J_0,N$.
\end{lemma}
\begin{proof}
For $x$ in one output cell, its \emph{entire} left domain complement
has $x-y\ge0$, and its entire right complement has $x-y\le0$.
Apply \eqref{eq:Pn} to each. On that output cell the ranges lie in
the spans of $N+1$ functions $u(x)e^{-\gamma x-2jx}$ on one side
and $N+1$ reflected exponentials on the other. The input factors are
square integrable: $v_I(y)\le e^{\alpha y/2}$ and the slowest decay
exponent is $1/2$. Summing output cells costs $2J_0(N+1)$ rank,
not a cost per pair of cells. This includes all inputs outside $I$.
The squared error is bounded using \eqref{eq:HS}, with $k_\pm$
replaced by the two exponential error envelopes. Thus its
Hilbert--Schmidt norm is at most $C3^{-N}$.

Within a cell, write $z=x-y\in(-\ell,\ell)$. The periodic images
$k_\pm(z+p\ell)$ have fixed positive sign of their argument for
$p\ge1$, and negative sign for $p\le-1$. Summing the separated
exponentials over $p$ is a geometric series and preserves the rank
$N+1$ on each side. The pointwise remainder is bounded by
\[
 C3^{-N}\{e^{-\gamma(\ell+z)}+e^{-\gamma'(\ell-z)}\},
 \qquad \gamma,\gamma'\in\{1/2,3/2\},
\]
with $C$ uniform for $\ell\ge1$. In one cell,
$e^{\alpha y}\le e^{\alpha\ell}e^{\alpha x}$; integrate in $y$
and sum over $x$. The bound \eqref{eq:HS} gives squared error at most
$C\ell e^{\alpha\ell}3^{-2N}$. The additional rank is
$2J_0(N+1)$, proving \eqref{eq:local}.
\end{proof}

Use $\ell=\sqrt L$, $J_0=O(L/\ell)$ and $N=\lceil C_1L\rceil$.
The norm error can then be made smaller than $e^{-C_2L}$ for any
fixed $C_2$ by increasing $C_1$, while the rank cost is $O(L^{3/2})$.
Apply the lemma in \eqref{eq:exactfactor} to write
$\cW=\cA+F+E$. Its uniform norm bound follows from \eqref{eq:HS}.
Set $\widetilde W=\cA+F=\cW-E$, which is also uniformly bounded.
Then
\begin{equation}\label{eq:producterror}
 \cW K\cW^*=\cA K\cA^*+R_L+Z_L,\qquad
 \rank R_L\le2\rank F=O(L^{3/2}),\qquad
 \|Z_L\|\le C e^{-C_2L}.
\end{equation}
Indeed
$\widetilde W K\widetilde W^*-\cA K\cA^*
 =FK\widetilde W^*+\cA KF^*$ has the stated rank. Bounding the
small error using $\widetilde W$, rather than $F$ separately, avoids
an uncontrolled norm from the finite-rank correction. For compact
operators the singular-value perturbation inequalities imply
\begin{equation}\label{eq:rankperturb}
 N_{\Delta f}(\varepsilon)\ge
 N_{\cA K\cA^*}(\varepsilon+\|Z_L\|)-\rank R_L,
\end{equation}
where here $N_T(\varepsilon)=\#\{j:\sigma_j(T)>\varepsilon\}$.

\section{The periodic channels and the sharp lower limits}

The Fourier transform convention is
$\widehat k(\xi)=\int_\R k(x)e^{-i\xi x}\,dx$.
The beta integral for $(2\cosh x)^{-1}$ and an allowed complex shift
give
\begin{equation}\label{eq:fourier}
 h_+(\xi)=\frac\pi2\sech\left(\frac{\pi\xi}2-\frac{i\pi}4\right),
 \qquad h_- =\overline{h_+},\qquad
 |h_\pm(\xi)|^2=\frac{\pi^2}{2}\sech(\pi\xi).
\end{equation}
On a circle of length $\ell$ these are the multipliers at
$\xi_k=2\pi k/\ell$, with no extra factor $\ell$.
Let $Q$ have positive multiplier $|h_+|$ and let
$U=\diag(U_+,U_-)$ contain their unitary phase multipliers.
Write $\widehat H=H_0/\sqrt2$, a constant unitary matrix. On one cell,
\begin{align}
 \cA&=\sqrt2 M_u U\widehat H\diag(QM_{v_I},QM_{v_I}),\nonumber\\
 \cA K\cA^*
 &=2M_u U\widehat H\diag(\kappa S,-\lambda S)
                 \widehat H^*U^*M_u,
 \qquad S=Q M_{v_I^2}Q\succeq (\inf v_I)^2 Q^2.\label{eq:cell}
\end{align}
The positive spatial weights $u,v_I$ and their reciprocals are bounded
on each finite cell.
The two central blocks have singular values $|\kappa|$ and $|\lambda|$
times the eigenvalues of the \emph{same positive} $S$. The unitary
factors preserve singular values, and two outer $M_u$ factors lower
them by at most $(\inf u)^2$. This proves a positive singular-value
lower for arbitrary complex coefficients. It requires neither
positivity of $K$ nor an indefinite Araki--Lieb--Thirring inequality.

Choose cells covering $[-A_0L,B_0L]$, with
$A_0>1/\alpha$, $B_0>1/(\beta+\delta)$, enlarged to cell endpoints.
Define
\[
 g(x)=\alpha(-x)_++(\beta+\delta)x_+.
\]
On the cell beginning at $x_j$, \eqref{eq:weights} and
$\omega_I\ge\omega_0$ give
\begin{equation}\label{eq:minweights}
 (\inf u)^2(\inf v_I)^2\ge c e^{-g(x_j)-C\ell}.
\end{equation}
On the left $u^2$ is constant; on the right it loses at most
$e^{-(1+\delta)\ell}$ across a cell. A cell crossing zero satisfies
the same bound after increasing $C$.

The two amplitudes must be retained separately when coefficients vary.
The resulting simultaneous profile is stronger than a lower bound at
one threshold and will serve both the moment and dimension limits.

\begin{lemma}[A uniform finite-matrix profile]\label{lem:uniformprofile}
Fix $\delta>0$ and $K_0>0$, and put
\[
 b_\delta=\frac1{2\pi^2}
       \left(\frac1\alpha+\frac1{\beta+\delta}\right)<b_\alpha.
\]
For every sufficiently large $L$ there is one finite Hermitian pair
$(A_L,B_L)$, independent of $\theta$, with rank-one increment of norm
$t_L\in(0,1]$, such that for all $|\theta|\le1$ and $0\le x\le K_0$,
\begin{equation}\label{eq:uniformprofile}
 \frac{N_{\Delta_{\theta,L}}(e^{-xL})}{L^2}
 \ge b_\delta\{x^2+(x-u_{\theta,L})_+^2\}
              -C_{\alpha,\delta,K_0}L^{-1/2},\qquad
 u_{\theta,L}=-\frac{\log|\theta|}{L},
\end{equation}
where $\Delta_{\theta,L}=f_\theta(B_L)-f_\theta(A_L)$.
The second term is zero if $\theta=0$.
\end{lemma}
\begin{proof}
Use the centered construction on cells covering
$[-(K_0+1)L/\alpha,(K_0+1)L/(\beta+\delta)]$.
The singular pencil has $K_\theta=\diag(1,-\theta)$ and
$\|K_\theta\|\le1$. Therefore the constants in
\eqref{eq:producterror} are uniform in $\theta$. Choose the localization
degree so that its norm error is below $C e^{-(K_0+3)L}$ and its rank
error is $O_{\alpha,\delta,K_0}(L^{3/2})$.

By \eqref{eq:fourier}--\eqref{eq:minweights}, the two channels on a
cell starting at $x_j$ have the ordered multiset singular lower
\[
 c a_\sigma e^{-g(x_j)-C\ell}e^{-2\pi^2|k|/\ell},
 \qquad k\in\mathbb Z,\qquad a_\sigma\in\{1,|\theta|\},
\]
with $c,C$ independent of $\theta$. In particular the count at
$4e^{-xL}$ is at least the sum over cells of
\begin{equation}\label{eq:cellcounts}
 \frac\ell{\pi^2}[xL-g(x_j)-C\ell-C_0]_+-1
 +\frac\ell{\pi^2}[xL+\log|\theta|-g(x_j)-C\ell-C_0]_+-1.
\end{equation}
For $\theta=0$ omit the second channel; retaining its $-1$ error only
weakens the lower. For any real $h$,
\[
 \int_\R[h-g(y)]_+\,dy
 =\frac{h_+^2}{2}
       \left(\frac1\alpha+\frac1{\beta+\delta}\right).
\]
Both effective positive depths in \eqref{eq:cellcounts} are at most
$K_0L$. The cell-minimum loss $O(\ell)$ therefore costs
$O_{K_0}(L\ell)$ in the squared phase volume, uniformly even when the
second depth tends to zero. The Riemann error has that order, and
the integer count losses are $O_{K_0}(L/\ell)$. With $\ell=\sqrt L$,
\begin{equation}\label{eq:modelprofile}
 N_{\cA K_\theta\cA^*}(4e^{-xL})
 \ge b_\delta L^2\{x^2+(x-u_{\theta,L})_+^2\}
                          -C_{\alpha,\delta,K_0}L^{3/2}
\end{equation}
uniformly over the stated $x,\theta$. The chosen interval contains
both triangular phase-volume supports.

Apply \eqref{eq:rankperturb}. The norm error is negligible relative
to $e^{-xL}$ uniformly on $[0,K_0]$, so the factor $4$ leaves a
strict threshold margin; the discarded rank costs $O(L^{3/2})$.
To transfer to finite matrices, approximate $D$ uniformly by finite
step functions, keeping $w$ fixed. At this fixed $L$, continuous
functional calculus lets us make \emph{both} $g_+$ and $g_-$ corrections
change by less than $e^{-(K_0+3)L}$. Thus every $f_\theta$ correction
changes by at most twice that amount, simultaneously for all
$|\theta|\le1$. The finite cell-vector spans and their two-copy sums
reduce both endpoints; the corrections vanish on the orthogonal
complements. These are genuine finite matrices. Their remaining
strict margin and \eqref{eq:modelprofile} prove
\eqref{eq:uniformprofile} for all $x,\theta$ with one finite pair.
The increment norm is still $t_L=\int_I\rho\in(0,1]$.
\end{proof}

\begin{proof}[Proof of Theorem~\ref{thm:uniform}]
Scalar homogeneity, a scalar phase, and the exchange
$g_+(-x)=g_-(x)$ reduce the lower to $f_\theta$, $|\theta|\le1$.
Put $L=1/s$ and $u=-s\log|\theta|$, including $u=\infty$.
For the same finite pair in Lemma~\ref{lem:uniformprofile}, the input
power denominator $t_L^{\alpha s}$ is at most one. Layer cake over
positive logarithmic thresholds gives
\begin{equation}\label{eq:uniformmomentlower}
 s^2\cH_s(f_\theta)\ge
 b_\delta\int_0^{K_0}e^{-x}\{x^2+(x-u)_+^2\}\,dx
                    -C_{\alpha,\delta,K_0}\sqrt s.
\end{equation}
Values above one only add a nonnegative contribution. The identity
\[
 \int_0^\infty e^{-x}\{x^2+(x-u)_+^2\}\,dx
        =2(1+e^{-u})=2(1+|\theta|^s)
\]
holds also for $u=\infty$. Since $(x-u)_+^2\le x^2$, the omitted
tail is at most $2T_{K_0}$, where
$T_{K_0}=\int_{K_0}^\infty e^{-x}x^2\,dx\to0$.
Divide \eqref{eq:uniformmomentlower} by $1+|\theta|^s\ge1$ to get
\[
 \frac{s^2\cH_s(f_\theta)}{1+|\theta|^s}
 \ge2b_\delta-2b_\delta T_{K_0}
                       -C_{\alpha,\delta,K_0}\sqrt s.
\]
First take $s\downarrow0$, then $K_0\to\infty$, then
$\delta\downarrow0$. The lower is uniform in $\theta$ and matches
\eqref{eq:coefficientupper}. This proves \eqref{eq:uniform} without
interchanging a supremum and an integral.
\end{proof}

\begin{proof}[Proof of Theorem~\ref{thm:accuracy} and
Corollary~\ref{thm:main}]
Set $K_0=x=1$ in Lemma~\ref{lem:uniformprofile}. Only after preserving
the counts, scale both endpoints by $1/t_L$; homogeneity increases
every output singular value and gives exact unit input norm.
Using $r$ direct copies gives the uniform lower
\[
 \cN_{f_\theta,r}(L)\ge
 r b_\delta\{L^2+[L+\log|\theta|]_+^2\}
                        -O_{\alpha,\delta,r}(L^{3/2}).
\]
Together with \eqref{eq:uniformupper}, let $L\to\infty$ first, then
$d\uparrow\pi/2$ and $\delta\downarrow0$.
This proves \eqref{eq:accuracy}. For fixed nonzero coefficients,
$|\kappa|^s+|\lambda|^s\to m$, giving the fixed moment limit.
For the count, a fixed scalar factor shifts $L$ by a constant;
a fixed nonzero smaller branch has $L^{-1}\log|\theta|\to0$.
This proves \eqref{eq:main}; the zero function is immediate.
\end{proof}

\section{A Gaussian model for arbitrary pole sets}

The compression argument needs both left and right column Grams,
including when quadrature poles are exact conjugates. Direct polynomial
identities give a single construction for every pole configuration.

\begin{lemma}[A $q$-state singular-spectrum realization]\label{lem:qstate}
Let $q\ge1$, let $z_1,\ldots,z_q$ be distinct nonreal poles, and let
$d_1,\ldots,d_q\in\C$. Exact conjugate pairs are allowed.
For $A=A^*$, $B=A+vv^*$, $\|v\|=1$, there are Hermitian
$A'$, $B'=A'+v'v'^*$ of size at most $q$, with
$0<\|v'\|^2\le1$, such that
\begin{equation}\label{eq:matchedquad}
 Q=\sum_{j=1}^q d_j[R_B(z_j)-R_A(z_j)],\qquad
 Q'=\sum_{j=1}^q d_j[R_{B'}(z_j)-R_{A'}(z_j)]
\end{equation}
have identical singular values after zero padding.
\end{lemma}
\begin{proof}
Let $\mu$ be the probability spectral measure of $(A,v)$, and put
\[
 m(z)=\int\frac{d\mu(t)}{t-z},\qquad
 w(t)=\prod_{j=1}^q(t-z_j)(t-\bar z_j)>0\quad(t\in\R).
\]
If $\mu$ has at most $q$ support points, use that measure exactly.
Otherwise take the ordinary positive $q$-node Gaussian rule for
$d\nu=d\mu/w$, exact on polynomials through degree $2q-1$.
Let its nodes be $t_k$, weights $a_k>0$, and monic orthogonal
polynomial $p_q$; its nodes are the zeros of $p_q$.
This is the classical modified-measure rational quadrature construction
\cite[Theorem 3.1 and Section 6(ii)]{Gautschi}.
Ordinary Gaussian exactness follows by dividing a polynomial of degree
$2q-1$ by $p_q$, using orthogonality for the quotient and interpolation
for the remainder. The weights are integrals of squared Lagrange
polynomials and are positive; the orthogonal polynomial has real
simple zeros inside the support interval.

Reweight by $b_k=a_kw(t_k)>0$ and let
$\mu'=\sum b_k\delta_{t_k}$. Since $w-p_q^2$ has degree at most
$2q-1$ and $p_q$ vanishes at the nodes,
\begin{equation}\label{eq:massdefect}
 1-\sum_k b_k=\int w\,d\nu-\sum_k a_kw(t_k)
                =\int p_q^2\,d\nu\ge0.
\end{equation}
The new mass is positive. Set $A'=\diag(t_k)$,
$v'=(\sqrt{b_k})$, and $B'=A'+v'v'^*$. Then
$0<\|v'\|^2=\sum b_k\le1$. In the finite-support case the mass is
one and the same conclusions hold.

The functions $w(t)/(t-z_j)$ and $w(t)/(t-\bar z_j)$ are polynomials
of degree $2q-1$, so $\mu'$ matches $m$ at all poles and conjugates.
Define columns $U_j=R_A(z_j)v$, $V_j=R_A(\bar z_j)v$. Their two
Grams are
\begin{equation}\label{eq:directGrams}
 (U^*U)_{ij}=\int\frac{d\mu(t)}{(t-\bar z_i)(t-z_j)},\qquad
 (V^*V)_{ij}=\int\frac{d\mu(t)}{(t-z_i)(t-\bar z_j)}.
\end{equation}
Multiply either integrand by $w$. The result is a polynomial of
degree $2q-2$, even when $z_j=\bar z_i$: in that case $w$ contains
the required repeated factor from the two indexed poles. Gaussian
exactness therefore reproduces \emph{both} Grams directly. No divided
difference with a zero denominator, and no separate resolvent-derivative
matching assumption, is used.

Woodbury gives
\begin{equation}\label{eq:Woodone}
 R_B(z_j)-R_A(z_j)=-\frac{U_jV_j^*}{1+m(z_j)}.
\end{equation}
The denominator is nonzero since $B$ is Hermitian and $z_j$ is nonreal;
its value is matched by $\mu'$. Thus $Q=UDV^*$ and $Q'=U'DV'^*$
have the same diagonal $D$ and the same two column Grams.
The Gram identities induce isometries on the left column spans and,
separately, on the right column spans, proving equality of singular
values. They need not define a single common Hermitian congruence.
\end{proof}

Taking $q+1$ Gaussian nodes would match mass exactly. The $q$-node
mass inequality already suffices for the common envelope and only
improves a H\"older ratio.
The following consequence uses that inequality before normalization.

\begin{corollary}[Uniform spectral compression]\label{cor:compression}
Fix $0<d<\pi/2$. For $f_\theta=g_++\theta g_-$ with $|\theta|\le1$,
a unit rank-one pair can be replaced by a pair of size
\begin{equation}\label{eq:compressionstates}
 n'\le b_d\{h^2+[h+\log|\theta|]_+^2\}
                             +O_{\alpha,d}(h^{7/4}+1)
\end{equation}
whose increment has norm in $(0,1]$ and whose correction satisfies,
after zero padding,
\begin{equation}\label{eq:spectralcompression}
 |\sigma_j(\Delta)-\sigma_j(\Delta')|\le2e^{-h}.
\end{equation}
The bounds are uniform in $\theta$ and the background as $h\to\infty$.
\end{corollary}
\begin{proof}
Use the separate branch quadratures from \eqref{eq:uniformupper} at
target depth $h$, choosing the fixed additive depth constants to make
the combined operator error at most $e^{-h}$. Merge repeated poles;
the total number $q$ satisfies \eqref{eq:compressionstates}.
Lemma~\ref{lem:qstate} gives a model in at most $q$ states with the
same quadrature singular values. Its smaller increment norm still
satisfies the same unit-input envelope, so the \emph{same} nodes and
weights approximate its actual correction to $e^{-h}$ as well.
The singular-value perturbation inequality on both sides gives
\eqref{eq:spectralcompression}. Keep $v'$ unnormalized during this
comparison. Its denominator $\|v'\|^{2\alpha s}\le1$ helps a moment
lower; to impose unit input for a count, scale both endpoints only
after preserving the desired singular values.
\end{proof}

These are exact existence and approximation statements. They give
neither numerical conditioning bounds for the Gaussian model nor a
runtime for shifted solves. The rational quadrature and matrix-update
principles are classical antecedents \cite{Gautschi,BCKS}.

\section{The joint dimension law and endpoint regimes}

\begin{proof}[Proof of Theorem~\ref{thm:joint}]
First prove the upper with $n_s s^2\to t\in(0,\infty)$.
For $b_d$ in \eqref{eq:bd}, set
$p_{u,d}(x)=b_d[x^2+(x-u)_+^2]$.
Rank-one reduction within dimension $n_s$,
\eqref{eq:uniformupper}, and the dimension cap give
\begin{equation}\label{eq:scaledcap}
 s^2N_\Delta(e^{-x/s})\le
 \min\{n_s s^2,p_{u_s,d}(x)
                  +C_{\alpha,d}(s^{1/4}x^{7/4}+s^2)\},\quad x\ge0.
\end{equation}
A common output norm bound is $M_\alpha\ge1$.
Values above one contribute at most $n_s(M_\alpha^s-1)$ to the
moment, whose product with $s^2$ vanishes. Layer cake and dominated
convergence in \eqref{eq:scaledcap} give the upper
$\int_0^\infty e^{-x}\min\{t,p_{u,d}(x)\}\,dx$.
The same reasoning applies for $u=\infty$ since the second term
vanishes at every fixed $x$. Then take $d\uparrow\pi/2$.

For the lower fix $\delta>0$, $0<d<\pi/2$, and $a_0>0$ with
\begin{equation}\label{eq:room}
 p_{u,d}(a_0)<t.
\end{equation}
Use Lemma~\ref{lem:uniformprofile} with $L=1/s$ and $K_0>a_0+1$.
Scale both original endpoints by $1/t_L$ to give unit input norm;
this increases all their output singular values. Apply
Corollary~\ref{cor:compression} at depth $h=a_0/s$. Its state count is
\[
 n'\le s^{-2}b_d\{a_0^2+(a_0-u_s)_+^2\}+o(s^{-2})<n_s
\]
eventually, by \eqref{eq:room}. The compressed increment has norm
at most one. If $0<x'<x<a_0$, then
\[
 e^{-x'/s}-2e^{-a_0/s}>e^{-x/s}
\]
for small $s$. Equations~\eqref{eq:uniformprofile} and
\eqref{eq:spectralcompression}, followed by $x'\uparrow x$, give
\begin{equation}\label{eq:compressedprofile}
 \liminf_{s\downarrow0}s^2N_{\Delta'_s}(e^{-x/s})
       \ge b_\delta\{x^2+(x-u)_+^2\},\qquad 0<x<a_0.
\end{equation}
For $x\ge a_0$, monotonicity and $x'\uparrow a_0$ supply the constant
lower $b_\delta\{a_0^2+(a_0-u)_+^2\}$.
The input power denominator is at most one. Fatou therefore gives
\begin{align*}
 \liminf_{s\downarrow0}s^2\cH_{s,n_s}(f_{\theta_s})
 &\ge b_\delta\int_0^\infty e^{-x}
 \{\min(x,a_0)^2+[\min(x,a_0)-u]_+^2\}\,dx.
\end{align*}
Let $\delta\downarrow0$, then increase $a_0$ to the solution of
$p_{u,d}(a_0)=t$, and finally take $d\uparrow\pi/2$.
The limiting profile is exactly $\min\{t,p_u(x)\}$, proving
\eqref{eq:joint}. This uses one actual family, rather than integrating
supremal counts from different matrices.

For the explicit formulas, integration by parts gives
$F(t,u)=\int_0^{a_0}e^{-x}p_u'(x)\,dx$.
If $a_0\le u$, only $2b_\alpha\int_0^{a_0}xe^{-x}\,dx$ occurs.
If $a_0>u$, add
$2b_\alpha e^{-u}\int_0^{a_0-u}ye^{-y}\,dy$.
Solving $p_u(a_0)=t$ gives \eqref{eq:jointone}--\eqref{eq:jointtwo},
which agree at $t=b_\alpha u^2$.

Uniformity on $t_0\le n s^2\le t_1$, for $0<t_0<t_1<\infty$,
follows from the sequential statement: a violating sequence has a
subsequence with $n s^2\to t$ and $u_s\to u$ in the compact extended
interval $[0,\infty]$. The function $F$ is continuous there.
In fact it is already independent of $u$ for
$u\ge\sqrt{t_1/b_\alpha}$, by \eqref{eq:jointone}.
For $n_s s^2\to\infty$, the universal upper in
Theorem~\ref{thm:uniform} gives $2b_\alpha(1+e^{-u})$.
Embed examples of size $\lfloor T/s^2\rfloor$, use
\eqref{eq:joint} at fixed $T$, and let $T\to\infty$ for the reverse.
The scalar normalization in \eqref{eq:jointscale} is exact by
homogeneity.
\end{proof}

The accuracy cap has the same two depth contributions.
\begin{corollary}\label{cor:jointcounts}
Let $r$ be a fixed positive integer, $n_L/L^2\to t\in(0,\infty)$,
$|\theta_L|\le1$, and $-L^{-1}\log|\theta_L|\to\tau\in[0,\infty]$.
Then
\begin{equation}\label{eq:jointcounts}
 \frac{\cN_{f_{\theta_L},r,n_L}(L)}{L^2}
 \longrightarrow\min\{t,r b_\alpha[1+(1-\tau)_+^2]\}.
\end{equation}
\end{corollary}
\begin{proof}
The dimension cap and \eqref{eq:uniformupper} give the upper.
For $r=1$, take the profile family at scale $L$, give it unit input
by scaling, and compress at depth $a_0L$ with $0<a_0<1$ and
$p_{\tau,d}(a_0)<t$. For every $0<x<a_0$, original singular values
above $e^{-xL}$ remain above $e^{-L}$ after the $O(e^{-a_0L})$ error.
There are at least
$[b_\delta\{x^2+(x-\tau)_+^2\}+o(1)]L^2$ of them.
Increase $x,a_0$ to the smaller of $1$ and the inverse of
$p_{\tau,d}$ at $t$, then take $\delta\downarrow0$ and
$d\uparrow\pi/2$. Normalize the compressed increment only after
count preservation if its norm is below one. This gives
$\min\{t,p_\tau(1)\}$. For fixed $r$, use $r$ direct copies with
block sizes at most $\lfloor n_L/r\rfloor$; their combined input
norm is one and their increment rank is $r$.
\end{proof}

\begin{proof}[Proof of Corollary~\ref{thm:dimension}]
For fixed $f\ne0$, scalar normalization gives $u_s\to0$ for two
active branches and $u_s=\infty$ for one. Theorem~\ref{thm:joint}
therefore gives \eqref{eq:transition} and its infinite-$t$ endpoint.
Corollary~\ref{cor:jointcounts} gives \eqref{eq:cappedcount}; fixed
scalar amplitudes shift $L$ by only a constant.

For fixed $n$, choose positive diagonal $A$ with distinct eigenvalues
and a unit $v$ with every coordinate nonzero. The PSD correction
$(A+vv^*)^\alpha-A^\alpha$ is strictly positive definite: its positive
resolvent integral can have zero quadratic form only if
$v^*(A+t)^{-1}z$ vanishes identically, forcing every residue and $z$
to vanish. If the coefficient $a\ne0$ in \eqref{eq:f}, this gives a
full-rank correction for $f$. If $a=0,b\ne0$, use the mirrored
negative endpoints. Its power trace tends to $n$, while rank-one
reduction gives $\cH_{s,n}(f)\le nM_f^s$.

For $n\to\infty$ and $n s^2\to0$, work at auxiliary scale
$L=\sqrt n$ with the fixed scalar-normalized function $f_\theta$.
Choose $d$ near $\pi/2$, $\delta$ small, and $0<x_0<a_0$ such that
$m b_d a_0^2<1$ and $m b_\delta x_0^2>1-\epsilon$.
This is possible since $b_d\downarrow b_\alpha$ and
$b_\delta\uparrow b_\alpha$. In this auxiliary scale,
$-L^{-1}\log|\theta|\to0$ if the second branch is present, and is
infinite otherwise. The uniform profile and compression at depth
$a_0\sqrt n$ therefore fit in $n$ states and preserve at least
$(1-\epsilon)n$ singular values above $e^{-C_\epsilon\sqrt n}$.
The compressed denominator is at most one, so the original exponent
$s$ gives a lower
$(1-\epsilon)n e^{-C_\epsilon s\sqrt n}$, up to the fixed scalar
factor to power $s$, which tends to one. This tends to
$(1-\epsilon)n$. The upper $nM_f^s$ matches. Bounded $n$ follows
from the fixed-size argument, completing all endpoints.
\end{proof}

\section{Bounded inputs in semifinite algebras}

This extension concerns a nonzero von Neumann algebra $(\mathcal M,\tau)$
with a normal semifinite faithful trace. Define $\cH_s^f(\mathcal M)$
by \eqref{eq:H}, with bounded Hermitian inputs in $\mathcal M$ and
$0<\tau(|B-A|^{\alpha s})<\infty$, replacing $\tr$ by $\tau$.
No assertion about all unbounded measurable inputs or type-III
Haagerup spaces is included. Let $D(\mathcal M)$ be its matrix degree:
the least $D$ for which it is $D$-subhomogeneous, or infinity.
At finite degree it is a finite direct sum of
$M_k(\mathcal A_k)$, $k\le D$, over abelian centers, with a nonzero
$D$ summand.

\begin{theorem}\label{thm:algebra}
For fixed nonzero $f$, if $D(\mathcal M)<\infty$, then
$\cH_s^f(\mathcal M)=\cH_{s,D(\mathcal M)}(f)$ and
$\cH_s^f(\mathcal M)\to D(\mathcal M)$. If $D(\mathcal M)=\infty$,
then $s^2\cH_s^f(\mathcal M)\to2c_f$.
\end{theorem}
\begin{proof}
For a finite-trace projection increment $E=p$, every pole difference
has support trace at most $\tau(p)$, by the resolvent identity and
polar decomposition. Thus the Gaussian quadrature has support trace
at most $m q_L\tau(p)$. The generalized singular-number approximation
principle \cite[Proposition 2.4]{FK} gives
\[
 \tau(1_{(e^{-L},\infty)}(|\Delta|))
 \le\tau(p)[c_d L^2+C_dL^{7/4}+K_d].
\]
The constant term also multiplies $\tau(p)$, since the bound at one
fixed threshold comes from a finite number of the same samples.
Layer cake gives a moment upper $U_{s,d}\tau(p)$ with
$\limsup s^2U_{s,d}\le2c_d$, uniformly in the algebra and the
background. Positive scaling handles $E=t p$; negative increments
are reversed positive ones.

The trace power triangle for $0<s\le1$
\cite[Proposition 4.6(i)]{FK} telescopes signed spectral step increments,
giving denominator $\sum |t_j|^{\alpha s}\tau(p_j)$.
For bounded general $E$, take finite spectral step approximants away
from zero, converging in norm, with their $\alpha s$ powers dominated
by $|E|^{\alpha s}$. Their supports have finite trace. Continuous
functional calculus and normal-trace Fatou give the same inequality
for $E$. The latter Fatou step follows by testing on finite-trace
projections and taking their supremum. Hence
$\limsup s^2\cH_s^f(\mathcal M)\le2c_f$ after moving $d$ to the
strip boundary, directly in $\mathcal M$.

If matrix degree is unbounded, for every $n$ there is an injective
nonzero homomorphism $M_n\to\mathcal M\cap L_1(\mathcal M,\tau)$;
the precise semifinite statement is \cite[Proposition 3.32]{Arhancet}.
Its trace is a positive finite multiple of matrix trace, which cancels
in the ratio. Since $f(0)=0$, extension from the corner by zero respects
functional calculus. Thus
$\cH_s^f(\mathcal M)\ge\sup_n\cH_{s,n}(f)$, proving the lower.
No hyperfiniteness assumption is used.

At finite degree, apply the $k$-by-$k$ inequality fiberwise and
integrate against the central trace to get the upper
$\cH_{s,D}(f)$. A finite-positive-trace central piece in the nonzero
$D$ summand and constant matrices give the reverse. The fixed-dimension
limit proved above completes the assertion.
\end{proof}

Thus the finite small-$s$ limit recovers bounded matrix degree, whereas
unbounded degree forces the full universal channel-dependent divergence.
This is a structural deduction from the sharp matrix theorem, not an
independent priority claim. The dimension transition describes precisely
when the extra channel becomes visible; for every fixed matrix size
the optimal power constants tend instead to that size.

\begin{thebibliography}{99}
\bibitem{Companion}
Working research manuscript.
\emph{Exact small-exponent constants for fractional matrix maps}.
November 2026. Companion manuscript,
\href{sharp-fractional-holder.pdf}{local PDF}.

\bibitem{Ricard}
\'E. Ricard. Fractional powers on noncommutative $L_p$ for $p<1$.
Adv. Math. 333 (2018), 194--211.
\url{https://arxiv.org/abs/1805.05677}.

\bibitem{HSZ}
J. Huang, F. Sukochev, D. Zanin.
Operator $\theta$-H\"older functions with respect to $\|\cdot\|_p$,
$0<p\le\infty$. J. London Math. Soc. (2022).
\url{https://arxiv.org/abs/2110.09708}.

\bibitem{BKS}
M. Sh. Birman, L. S. Koplienko, M. Z. Solomyak.
Estimates of the spectrum of a difference of fractional powers of
selfadjoint operators. Izv. Vyssh. Uchebn. Zaved. Mat., no. 3 (1975), 3--10.
\url{https://www.mathnet.ru/eng/ivm6330}.

\bibitem{Pekarskii}
A. A. Pekarskii. Best uniform rational approximations of Markov functions.
Algebra i Analiz 7(2) (1995), 121--132; St. Petersburg Math. J.
7(2) (1996), 277--285. \url{https://www.mathnet.ru/eng/aa551}.

\bibitem{Laptev}
A. A. Laptev. Spectral asymptotic behavior of a class of integral operators.
Math. Notes 16(5) (1974), 1038--1043; Russian original,
Mat. Zametki 16(5) (1974), 741--750.
\url{https://doi.org/10.1007/BF01149794}.

\bibitem{AZ}
J. Arazy, L. Zelenko. Directional operator differentiability of non-smooth
functions. J. Operator Theory 55(1) (2006), 49--90.
\url{https://jot.theta.ro/jot/archive/2006-055-001/2006-055-001-003.pdf}.

\bibitem{Stahl}
H. Stahl. Best uniform rational approximation of $x^\alpha$ on $[0,1]$.
Bull. Amer. Math. Soc. 28 (1993), 116--122.
\url{https://arxiv.org/abs/math/9301217}.

\bibitem{Mardvilko}
T. S. Mardvilko. Uniform rational approximation of the odd and even
Cauchy transforms. Sbornik: Mathematics 216(2) (2025), 239--256.
\url{https://doi.org/10.4213/sm10116e}.

\bibitem{BCKS}
B. Beckermann, A. Cortinovis, D. Kressner, M. Schweitzer.
Low-rank updates of matrix functions II: Rational Krylov methods.
SIAM J. Numer. Anal. 59 (2021), 1325--1347.
\url{https://arxiv.org/abs/2008.11501}.

\bibitem{Gautschi}
W. Gautschi. The use of rational functions in numerical quadrature.
J. Comput. Appl. Math. 133 (2001), 111--126.
\url{https://doi.org/10.1016/S0377-0427(00)00637-3}.

\bibitem{FK}
T. Fack, H. Kosaki. Generalized $s$-numbers of $\tau$-measurable operators.
Pacific J. Math. 123 (1986), 269--300.
\url{https://doi.org/10.2140/pjm.1986.123.269}.

\bibitem{Arhancet}
C. Arhancet. Quantum information theory via Fourier multipliers on
quantum groups. arXiv:2008.12019v9, Section 3.4 (2021).
\url{https://arxiv.org/html/2008.12019v9}.
\end{thebibliography}
\end{document}
